\documentclass[11pt]{article}
\usepackage[T1]{fontenc}
\usepackage{lmodern}
\usepackage{amsmath,amssymb,amsthm,mathtools}
\usepackage{array,booktabs,tabularx,longtable}
\usepackage{xcolor,geometry,microtype}
\usepackage{hyperref,cleveref}
\geometry{margin=1in}
\newtheorem{theorem}{Theorem}[section]
\title{A Common Article Workload}
\author{aiTeX qualification}
\date{September 2001}

\begin{document}
\maketitle
\tableofcontents

\section{Text and mathematics}\label{sec:text}
Office traffic and affine geometry exercise common ligatures. Accented text:
R\'esum\'e, d\'ej\`a vu, and Stra\ss e.

\begin{theorem}\label{thm:sum}
For every positive integer $n$,
\[
  \sum_{k=1}^{n} k = \frac{n(n+1)}{2}.
\]
\end{theorem}

\begin{align}
  (a+b)^2 &= a^2+2ab+b^2, \\
  \int_0^1 x^2\,dx &= \frac13.
\end{align}

\section{Tables and links}
The result in \cref{thm:sum} was introduced in \cref{sec:text}. Visit the
\href{https://example.invalid/aitex/matrix}{benchmark destination}.

\begin{tabularx}{\linewidth}{@{}l>{\raggedleft\arraybackslash}Xr@{}}
\toprule
Item & Description & Value \\
\midrule
Alpha & A deterministic row with flexible width & 13 \\
Beta & A second row in blue & \textcolor{blue}{21} \\
\bottomrule
\end{tabularx}

\section{Lists}
\begin{enumerate}
  \item A numbered item with \textbf{bold} and \emph{emphasized} text.
  \item A nested list:
    \begin{itemize}
      \item first nested point;
      \item second nested point.
    \end{itemize}
\end{enumerate}
\end{document}
